\documentclass[11pt]{article}
\usepackage{amsmath,amssymb}
\usepackage{xurl}
\usepackage[hidelinks]{hyperref}
% Shared commands for notes in docs/paper-gaps/.
%
% Two halves.  The link commands below are project identity: OWNER/REPO is the
% placeholder that scripts/bootstrap_project.py replaces with this project's
% GitHub slug and Pages site.  Everything after them is NOTATION, and notation
% belongs to the source paper: the definitions here are an example set, and a
% project replaces them with the macros of its own paper's style file.  The one
% rule that never changes: a command defined here must mean exactly what the
% source paper's macro means, or a note silently says something else.

\newcommand{\Lean}{\textsc{Lean}}
\newcommand{\leanid}[1]{\textsf{#1}}
\newcommand{\ghissue}[1]{\href{https://github.com/OWNER/REPO/issues/#1}{GitHub issue \##1}}
\newcommand{\ghpr}[1]{\href{https://github.com/OWNER/REPO/pull/#1}{PR \##1}}
% Local issue tracker (issues/NNNN-*.md in this repository); no remote URL.
\newcommand{\localissue}[1]{issue \##1 (\path{issues/})}
\newcommand{\doclink}[1]{\href{https://OWNER.github.io/REPO/docs/#1.html}{\path{#1}}}

% --- Notation: replace with the source paper's own macros --------------------
% \F, \N, \C, \R, \tr, \ot, \eps, \E, \bx, \bu, \bv, \by

\newcommand{\F}{\mathbb{F}}
\newcommand{\N}{\mathbb{N}}
\newcommand{\C}{\mathbb{C}}
\newcommand{\R}{\mathbb{R}}
\newcommand{\tr}{\mathrm{tr}}
\newcommand{\eps}{\epsilon}
\newcommand{\ot}{\otimes}
\newcommand{\E}{\mathop{\bf E\/}}

% bold random variables
\newcommand{\bu}{\boldsymbol{u}}
\newcommand{\bv}{\boldsymbol{v}}
\newcommand{\bx}{{\boldsymbol{x}}}
\newcommand{\by}{\boldsymbol{y}}

% T1 so literal < and > in \texttt placeholders typeset as themselves
\usepackage[T1]{fontenc}

% bra-ket
\usepackage{braket}

% --- Example structured spaces ---
\newcommand{\polyfunc}[3]{\mathcal{P}(#1, #2, #3)}
\newcommand{\polymeas}[3]{\mathrm{PolyMeas}(#1, #2, #3)}
\newcommand{\polysub}[3]{\mathrm{PolySub}(#1, #2, #3)}

% Approximation relations (paper uses \approx_\eps and \simeq_\zeta)
\newcommand{\approxeps}{\approx_{\varepsilon}}
\newcommand{\simgeq}{\simeq_{\zeta}}


\title{Prime-Characteristic and Binary Pauli Observables}
\author{}
\date{2026-09-02}

\begin{document}
\maketitle
\section{At a glance}
\begin{itemize}
  \item \textbf{Difficulty.} Medium.  The source-facing statements must retain
  arbitrary prime characteristic while the later Pauli basis test uses a
  separately named characteristic-two specialization.
  \item \textbf{Estimated weight.} The completed extension is the 513-line
  \path{MIPStarRE/QPBT/Algebra/PrimePauliBasis.lean} module.  It formalizes the
  source-general eigenvectors, rank-one projectors, eigenvalue calculation,
  projective-measurement facts, and both Fourier identities; the separately
  named binary consumer layer is unchanged.
  \item \textbf{Mathlib/project split.} Mathlib supplies finite matrix algebra
  and finite-field facts; the generalized Pauli definitions and their
  characteristic-two specialization are project-local.
  \item \textbf{Key mathematical inputs.} Unitarity of shift and phase
  matrices, the $p$-th-power identity, the Pauli eigenbasis expansions, and
  additive-character orthogonality for Fourier inversion.
  The scope correction is tracked in
  \href{https://github.com/Dengnifer/MIPStarRE-A/issues/16}{GitHub issue \#16}
  and descends from \localissue{0002}.
\end{itemize}

\section{Key theorem forms}
Let $K$ be a finite extension of $\mathbb F_p$, let
$\chi(x)=\exp(2\pi i\operatorname{Tr}_{K/\mathbb F_p}(x)/p)$, and let
$a,a',b\in K^n$.  The relevant identities are
\[
 \tau^W(a)\tau^W(a')=\tau^W(a+a'),\qquad
 \tau^W(a)^p=I,
\]
and
\[
 \tau^X(a)\tau^Z(b)=\chi(-a\mathbin{\cdot}b)
 \tau^Z(b)\tau^X(a).
\]
The explicit shift and phase matrices are unitary. The second identity implies
that every eigenvalue $\lambda$ of $\tau^W(a)$ satisfies $\lambda^p=1$. If
$V\leq K^n$ and $v\notin V^\perp$, Fourier cancellation takes the form
\[
 \mathbb E_{u\in V}\chi(u\mathbin{\cdot}v)=0.
\]
Writing $P_b^W$ for the rank-one projector onto the $b$-labelled vector in the
$W$ eigenbasis, the forward expansion and its Fourier inverse are
\[
 \tau^W(a)=\sum_{b\in K^n}\chi(a\mathbin{\cdot}b)P_b^W,
 \qquad
 P_b^W=\mathbb E_{a\in K^n}\chi(-a\mathbin{\cdot}b)\tau^W(a).
\]

\section{Source and formal scope}
The source defines the shift and phase matrices and states their product,
$p$-th-power, unitarity, spectral, and twisted-commutation properties in
\cite[lines 1056--1095]{Ji2020MIPStar}. It uses the same character for Fourier
cancellation over a subspace at lines 1124--1132.  The single-qudit forward and
inverse expansions occur at lines 1118--1139, and their multi-qudit forms occur
at lines 1151--1161.  These assertions use the characteristic-$p$ character;
they are not intrinsically binary.

The formal definitions \path{primeTauShift}, \path{primeTauPhase}, and
\path{primeTauObservable} use the canonical character \path{ffChar} in
characteristic $p$.  Their theorem statements cover the product, power,
unitarity, spectral, commutation, and cancellation identities for arbitrary
$p$.  The declarations \path{primePauliVec}, \path{primePauliFourier}, and
\path{primePauliProj} give the common eigenvectors, the Fourier matrix, and the
rank-one projectors of the same source definition.  In particular, the $X$
vector uses the negative phase $\chi(-e\cdot x)$, while
\path{primeTauObservable_mulVec_primePauliVec} proves the positive-phase
eigenvalue $\chi(a\cdot e)$.  Orthonormality, projectivity, positivity, and
completeness are proved by the accompanying declarations in
\path{MIPStarRE/QPBT/Algebra/PrimePauliBasis.lean}.

The source assumptions are that $p$ is prime, $K=\F_q$ is a finite extension
of $\F_p$, and the many-qudit carrier has positive finite dimension.  Lean
encodes these by \path{Fact p.Prime}, a finite field $K$ with an algebra from
\path{ZMod p}, and a finite decidable coordinate type $\iota$; the two
source-labelled Fourier identities additionally assume \path{Nonempty}, which
is exactly the condition $n\geq 1$.  The Boolean basis parameter uses
\path{true} for $X$ and \path{false} for $Z$.  These are boundary encodings of
the source domain, not additional mathematical hypotheses.

The declarations whose names begin with \path{tauObservable} or use
\path{phaseSign} remain the separately named specialization to $p=2$, where
$\chi(x)=(-1)^{\operatorname{Tr}(x)}$.  This specialization is appropriate for
the subsequent test, whose admissible sizes satisfy $q=2^k$; it is not
substituted for the general statements.

In characteristic two, $-x=x$, so the negative phase in the inverse expansion
is the same \path{phaseSign} used in the forward expansion.  Thus
\leanid{tauObservable\_eq\_sum\_pauliProj} and
\leanid{pauliProj\_eq\_avg\_tauObservable} prove precisely the two displayed
formulas for finite fields carrying an algebra over $\mathbb F_2$.  Their proof
uses additive-character orthogonality after transporting the coordinate type
along an equivalence with a finite ordinal; the restriction enters only through
the identification of the canonical character with \path{phaseSign}.

The proofs can be obtained directly from the displayed matrices.  A shift is a
permutation matrix, a phase matrix has diagonal entries of modulus one, and a
tensor product of such matrices is unitary.  Iterating the product identity
gives $\tau^W(a)^p=I$.  Applying this equality to a nonzero eigenvector gives
$(\lambda^p-1)v=0$, hence $\lambda^p=1$.  Finally, when
$v\notin V^\perp$, the map $u\mapsto\chi(u\cdot v)$ is a nontrivial character
of the finite additive group $V$, so its average vanishes.

\section{Blueprint and formalization}
The source-general definitions and results are recorded in
\path{blueprint/src/chapter/ch11_qpbt_algebra.tex} at
\path{def:generalized-pauli}, \path{lem:twisted-commutation}, and
\path{lem:pauli-observable-expansion}. Their Lean declarations are \leanid{primeTauShift},
\leanid{primeTauPhase}, \leanid{primeTauObservable}, and the accompanying
product, power, unitarity, spectral, commutation, and cancellation theorems,
together with \leanid{primePauliVec}, \leanid{primePauliFourier},
\leanid{primePauliProj}, and their orthonormality, eigenvalue, and projective
measurement theorems.  The source-general forward expansion and inverse are
\leanid{primeTauObservable\_eq\_sum\_primePauliProj} and
\leanid{primePauliProj\_eq\_avg\_primeTauObservable}.  They quantify an
arbitrary prime characteristic and use the corresponding additive character.

The separately named binary declarations are used by the later QPBT chapters.
In particular, \leanid{tauObservable\_eq\_sum\_pauliProj} is the forward
observable expansion and \leanid{pauliProj\_eq\_avg\_tauObservable} is its
normalized Fourier inverse.  They specialize the source formulas to
characteristic two and are not linked as implementations of the unrestricted
source node. Thus the blueprint, the general Lean interface, and the paper
agree; the binary layer is a consumer specialization rather than a restriction
of the source statement.

\section{Conclusion}
The source-facing shift, phase, observable, eigenvector, Fourier-transform, and
projector definitions all retain arbitrary prime characteristic and the
corresponding canonical character.  Their orthonormality, projective
measurement properties, eigenvalue identities, forward expansion, and Fourier
inverse are proved from the explicit matrices and additive-character
orthogonality.  The characteristic-two declarations remain separate
specializations for the QPBT parameter regime, and no binary definition or
consumer is changed by the source-general completion.  The corresponding
register row therefore proposes the terminal disposition \path{no-difference};
independent source-fidelity review of issue~\#673 and main's merge gate decide
whether that proposal is adopted.
\bibliographystyle{alpha}
\bibliography{references}
\end{document}
